
% ============================================================================
\clearpage
\section*{v4.10.0-r60 Source-Only: Triality, the Family Symmetry and Its Breaking to \texorpdfstring{\(A_4\)}{A4} (k275--k277, H127)}
\addcontentsline{toc}{section}{v4.10.0-r60 Source-Only: Triality, the Family Symmetry and Its Breaking to \texorpdfstring{\(A_4\)}{A4} (k275--k277, H127)}
\label{sec:v4100-r60-triality-a4}
% ============================================================================

\paragraph{Purpose.}
Layer r59.13 left one candidate open: a tower level on which \(\operatorname{Im}\mathbb H\) acts as a vector on matter, so that three generations
carry spin 1 rather than spin \(\tfrac12\); triality was named but not computed.  On 24 November 2026 one exploration (k275), one test with an
advance note (H127, k276) and one further exploration (k277) followed.  They were recorded as research notes with scripts and outputs.  This
layer collects them, each with its status as recorded at the time, on the author's decision to bring the Companion up to date before further
computation (the articles follow later).  It adds no new computation; every statement below is checked in the cited scripts.  k275 and k277 are
explorations and \emph{not blind}; H127 was fixed in an advance note before the computation.  Earlier layers are not rewritten.

\noindent\textbf{Source-only status.}
This v4.10.0-r60 layer is source-only.  It does not build a PDF, does not assign a new DOI,
and does not authorize a public release.  The v5.4.0 Version DOI
\href{https://doi.org/10.5281/zenodo.24500009}{10.5281/zenodo.24500009} and the Concept DOI
\href{https://doi.org/10.5281/zenodo.19210728}{10.5281/zenodo.19210728} are unchanged by it.

\subsection*{v4.10.0-r60 setting}
\label{subsec:v4100-r60-setting}

The carrier is the Albert algebra \(H_3(\mathbb O)\) of Hermitian \(3\times3\) octonion matrices with the Jordan product \(X\circ Y=\tfrac12(XY+YX)\);
its three off-diagonal octonions \(x_1=X_{23}\), \(x_2=X_{31}\), \(x_3=X_{12}\) are the three slots, and \(\operatorname{Aut}=F_4\supset\operatorname{Spin}(8)\) carries
triality.  Colour is the diagonal \(SU(3)_c\) (stabiliser of \(e_7\) in \(G_2\)), charge the slot-wise left multiplication
\(C=L_{e_7}\) (with \(R_{e_7}\), as in r35/r42).  ``Colour triplets'' are the \(e_1,\dots,e_6\) components (18 real dimensions), ``colour singlets'' the
diagonal and the \(1,e_7\) components (9), which form \(H_3(\mathbb C)\) with \(\mathbb C=\operatorname{span}(1,e_7)\).  Reading the three slots as
generations is a reading, not derived.  All metrics are the trace form \(\operatorname{Tr}(X\circ Y)\).  Frame convention F holds throughout.

\subsection*{v4.10.0-r60 results}
\label{subsec:v4100-r60-results}

\paragraph{Theoremlet v4.10.0-r60.1 (the family \(SO(3)\) from triality; k275, not blind).}
(a) \(\dim\operatorname{Der}H_3(\mathbb O)=52\) (\(f_4\)); \(su(3)_c\) and \(g_2\) lie in \(f_4\).  The commutant of \(su(3)_c\) in \(f_4\) is \(su(3)'\) (dimension 8),
the commutant of \(g_2\) is the real \(so(3)\) acting by \(X\mapsto RXR^{\mathsf T}\), contained in \(su(3)'\).
(b) Under this \(so(3)\) every colour triplet is spin 1 (six vectors); the colour singlets decompose as \(1\oplus3\oplus5\): the \(e_7\) parts are the
vector, diagonal and real parts spin \(0\oplus2\) (Wolfram exact).  The \(\varepsilon\) is built in, because the antisymmetric part of a real \(3\times3\)
matrix is a vector through \(\varepsilon\).
(c) \([so(3),L_{e_7}]\) vanishes exactly on the colour triplets and equals 2 on the colour singlets, where \(L_{e_7}\) mixes the real unit
(spin 2/0) with \(e_7\) (spin 1).  The part of \(su(3)'\) commuting with charge is abelian of dimension 2 (slot phases, no generation change).
(d) Of the 24 signed permutations with determinant \(+1\) exactly 12 commute with charge, all with even slot permutation; they form \(A_4\).  The odd
ones exchange \(x_{ij}\) with \(\bar x_{ij}\), and conjugation turns \(L_{e_7}\) into \(-R_{e_7}\).  Since the residual Lie algebra is zero and every finite
subgroup of \(SO(3)\) containing \(A_4\) is \(A_4\), \(S_4\) or \(A_5\) (an icosahedral rotation breaks), the remnant is exactly \(A_4\), with the slots in its
three-dimensional irreducible representation.
Thus the gap ``spin \(\tfrac12\) against spin 1'' of r59.13 is closed for the quarks; for the leptons it is located in \(\operatorname{span}(1,e_7)\).  The
\(A_4\) that r33 and r55 (H39) set as a slot symmetry is here the remnant of the family \(SO(3)\) under charge.  Literature: \(F_4\supset G_2\times SO(3)\),
\(F_4\supset SU(3)\times SU(3)\) (Dynkin; Dubois-Violette--Todorov; Boyle).  (\texttt{k275}, \texttt{k275b})

\paragraph{Theoremlet v4.10.0-r60.2 (H127: the charge breaking is pure spin 3; test).}
Advance note committed before the computation (\texttt{ba473c5}, 24 November 2026, 15:59 UTC).  Let \(P_s,P_t\) project onto singlets and triplets
and \(A_0=P_sCP_s\in\operatorname{End}(H_3(\mathbb C))\).  Disclosed in advance: on \(H_3(\mathbb C)\) one has \(L_{e_7}=R_{e_7}\), so the odd elements of
\(S_4\) send \(A_0\) to \(-A_0\); nonzero components can sit only at spins \(l\in\{3,6,7,9,10\}\).
\emph{Stage A (confirmed).}  The prediction was that \(l=3\) is present and leading.  Result: \(\operatorname{End}(1\oplus3\oplus5)\) decomposes as
\(3\times l_0+6\times l_1+6\times l_2+3\times l_3+1\times l_4\), and \(A_0\) lies \emph{entirely} in the spin-3 part, \(w_3=1\) and \(w_{l\neq3}=0\) (Wolfram exact,
\(\|A_0\|^2=6\)).  The breaking \(SO(3)\to A_4\) is a single spin-3 tensor of type \(xyz\) (\(t_{ijk}=|\varepsilon_{ijk}|\)), the object that \(A_4\)
flavon models set by hand.
\emph{Stage B (refuted by the wording; open in content).}  With the coupling \(K(A)=\sum_aL_{Y_a}AL_{Y_a}\) over a trace-orthonormal basis the effective
triplet operator \(P_tK^n(A_0)P_t\) vanishes for \(n=1,2,3\).  Post-hoc check: on \(\Lambda^2(27)=52\oplus273\) the map \(K\) is \(\tfrac32\) times the
projector onto \(f_4\) (rank 52, \(K(T)=1.5\,T\) for \(T\in so(3)\)), and \(C\) is orthogonal to \(f_4\).  The outcome follows from the choice of \(K\); the
reading ``quarks keep \(SO(3)\) effectively exactly'' given in the advance note is not covered by it.  This is recorded as a defect of the advance
note.  (\texttt{k276}, \texttt{k276b})

\paragraph{Theoremlet v4.10.0-r60.3 (no-go for the quarks; k277, not blind).}
On the colour triplets the \(A_4\)-invariant and the \(SO(3)\)-invariant operators coincide (both of dimension 36; every \(A_4\)-invariant commutes
with \(so(3)\) to \(10^{-15}\)).  Reason: spin \(1\otimes1\) has \(l=0,1,2\) only, and \(A_4\) has no invariant at \(l=1\) or \(l=2\).  Since \(C\) is
\(A_4\)-invariant on all 27 dimensions, every \(SO(3)\)-equivariant construction from \(C\), at any order and with any \(SO(3)\)-compatible coupling, is
\(SO(3)\)-invariant on the quarks.  With exact \(A_4\) the quarks therefore have degenerate masses and \(V_{\rm CKM}=\mathbb 1\): the leading-order
statement H39d is derived, not read.  Quark mixing requires that \(A_4\) itself breaks.  Stage B of r60.2 was empty for every coupling, not only for
the chosen one.  (\texttt{k277})

\paragraph{Theoremlet v4.10.0-r60.4 (the cubic term is the Albert norm; k277, not blind).}
Place a flavon \(v\) in the \(e_7\) parts of the three slots (there it is spin 1): \(X(v)\) is imaginary antisymmetric, \(\det X=0\),
\(\operatorname{Tr}X^2=2|v|^2\).  Its charge rotation \(Y=C\,X(v)\) is real symmetric with zero diagonal and entries \(-v_k\), and
\[
N(Y)=\det Y=-2\,v_1v_2v_3,\qquad \operatorname{Tr}Y=0,\qquad \operatorname{Tr}Y^2=2|v|^2
\]
(Wolfram exact).  The \(A_4\) cubic invariant \(xyz\) is the cubic norm of the Albert algebra after the charge rotation.
Mixed invariants of \(X\) and \(Y\) bring in \(\sum v_i^4\): \(\operatorname{Tr}((X\circ Y)^2)=2\sum_{i<j}v_i^2v_j^2\), \(\operatorname{Tr}(X^2\circ Y^2)=2\sum v_i^4\),
while \(\operatorname{Tr}((Y^2)^2)=2|v|^4\).  (\texttt{k277}, \texttt{k277b})

\paragraph{Theoremlet v4.10.0-r60.5 (vacuum alignment; k277, not blind).}
(a) A potential built only from \(Y\) depends only on \(|v|^2\) and \(xyz\) (the invariants of one element of a cubic Jordan algebra are \(\operatorname{Tr}\),
\(\operatorname{Tr}(\cdot)^2\), \(N\)).  The extrema of \(xyz\) on the sphere lie exactly at \(|v_i|=1/\sqrt3\) (8 solutions, Wolfram), so as soon as the
potential depends on \(xyz\) the minimum is of type \((1,1,1)\) with residual group \(\mathbb Z_3\), the reading \(G_D=\mathbb Z_3\) of the Dirac sectors in
H39.  Numerical minimisation agrees for \(\kappa=\pm0.5,0.05\).
(b) With mixed terms the potential is the general \(A_4\) potential \(-\mu^2|v|^2+\kappa\,xyz+\lambda|v|^4+\lambda'\sum v_i^4\); \(\lambda'=+0.3\) gives
\(\mathbb Z_3\), \(\lambda'=-0.15\) gives \((1,0,0)\) with residual \(\mathbb Z_2\) (the neutrino reading).  This is the known phase diagram (Altarelli--Feruglio);
the coefficients are free.
Hence no blind test of the alignment is possible at present; it needs a principle that fixes the coefficients (a candidate is the tower dynamics
of r59.14 with the cubic norm as natural third-order term, not computed).  The number H128 stays unassigned.

\paragraph{Remark v4.10.0-r60.6 (balance of the settings).}
Derived in this layer: the family \(SO(3)\) (commutant of \(G_2\) in \(F_4\)), its breaking to exactly \(A_4\) by charge, the spin-3 form of that breaking,
the quark no-go under exact \(A_4\) and the cubic invariant as Albert norm.  Set remain: the slots as generations, the flavon in the \(e_7\) parts,
the potential coefficients, every mass and every scale.  The open choice has moved from ``which slot symmetry'' (r33, r55) to ``which alignment'',
i.e.\ to the coefficients of a potential.  No item derives \(\beta\), a mass, a mixing angle or \(\Lambda\).

\subsection*{v4.10.0-r60 frames}
\label{subsec:v4100-r60-frames}

\begin{center}
\small
\begin{tabular}{@{}lll@{}}
\toprule
item & frame & comparison\\
\midrule
r60.1--r60.5 (Albert algebra, charge, flavon) & M (field value, matter readout) & none across frames\\
H39 readings (\(\mathbb Z_3\), \(\mathbb Z_2\)) & M & only as reading, no new map\\
T, P & not used & ---\\
\bottomrule
\end{tabular}
\end{center}

\subsection*{v4.10.0-r60 register entries}
\label{subsec:v4100-r60-register}

\begin{center}
\small
\begin{tabular}{@{}>{\raggedright\arraybackslash}p{0.40\textwidth}>{\raggedright\arraybackslash}p{0.52\textwidth}@{}}
\toprule
statement & status after v4.10.0-r60\\
\midrule
r60.1, r60.3--r60.5 & explorations, not blind; recorded\\
H127 stage A (charge breaking is spin 3) & confirmed\\
H127 stage B (leak order into the quarks) & refuted by the wording; empty for every \(SO(3)\)-equivariant coupling (r60.3)\\
\(A_4\) as slot symmetry (r33, r55) & derived as remnant of the family \(SO(3)\) under charge\\
H39d (\(V_{\rm CKM}=\mathbb 1\) at leading order) & derived under exact \(A_4\)\\
vacuum alignment \(\mathbb Z_3\) / \(\mathbb Z_2\) & coefficient dependent; no blind test (H128 unassigned)\\
OP1, OP3, OP4, OP5, Q1a & open\\
\bottomrule
\end{tabular}
\end{center}

\subsection*{v4.10.0-r60 notes}
\label{subsec:v4100-r60-notes}

{\raggedright
\texttt{notes/\allowbreak research/\allowbreak hilbert\_\allowbreak shift\_\allowbreak felder\_\allowbreak 2026\_\allowbreak 11\_\allowbreak 24/\allowbreak rueckwaerts/} with subfolders
\texttt{trialitaet/} (k275, k275b), \texttt{h127/} (k276, k276b) and \texttt{k277/} (k277, k277b);
\texttt{predictions/\allowbreak Vorab\_\allowbreak Notiz\_\allowbreak H127\_\allowbreak Ladungsbrechung\_\allowbreak Spin.md} and
\texttt{predictions/\allowbreak Ergebnis\_\allowbreak H127\_\allowbreak Ladungsbrechung\_\allowbreak Spin.md}.  Each folder holds scripts and their actual outputs.
\par}

\subsection*{v4.10.0-r60 ledger and guard update}
\label{subsec:v4100-r60-ledger-guard}

\[
R_{\rm chk}^{v4100\text{-}r60}=(1,1,0)
\]
for (status recorded; computations checked in the cited scripts; no new computation in this layer; generations, alignment and scales not derived).

The layer spends no guard:
\[
R_{\rm glob}^{(v4100\text{-}r60,+)}=(0,1,1,1,1),
\qquad
\Gamma_{\rm guard}^{v4100\text{-}r60}=(1,1,1,1,1),
\qquad
\beta_4=1/256\ \text{unspent}.
\]
The layer does not imply exact finite \(C5\) zero closure, exact all-readout terminality, \(C6/T2\), a global photon-clock/RH theorem,
or global all-readout closure.  OP1--OP5, C5 and C6/T2 remain open.
